\documentclass[a4paper,11pt]{article}

\usepackage[margin=1.2cm]{geometry}
\usepackage{tikz}
\usepackage[most]{tcolorbox}
\usepackage{array}
\usepackage{multicol}
\usepackage{enumitem}
\usepackage{titlesec}
\usepackage{xcolor}
\usepackage{graphicx}
\usepackage{amsmath}
\usepackage{tabularx}
\usepackage{lmodern}
\usepackage[T1]{fontenc}
\usepackage[utf8]{inputenc}
\pagestyle{empty}

\definecolor{mainblue}{HTML}{163A70}
\definecolor{lightblue}{HTML}{DCEBFF}
\definecolor{greenbox}{HTML}{DFF5E1}
\definecolor{purplebox}{HTML}{EEE3FF}
\definecolor{orangebox}{HTML}{FFF1DD}
\definecolor{cyanbox}{HTML}{DDF7F7}

\newtcolorbox{sectionbox}[2][]{
colback=#1,
colframe=mainblue,
fonttitle=\bfseries\large,
title=#2,
arc=4mm,
boxrule=1pt
}

\begin{document}

%==================== TITRE ====================

\begin{center}
\begin{tcolorbox}[
colback=mainblue,
colframe=mainblue,
width=\textwidth,
arc=3mm,
boxrule=0pt
]
\centering
{\Huge \textcolor{white}{\textbf{NOMENCLATURE DES ALCANES}}}\\[0.3cm]
{\Large \textcolor{white}{Règles simplifiées de nomenclature organique}}
\end{tcolorbox}
\end{center}

\vspace{0.4cm}
%==================== LIGNE 1 ====================
\begin{sectionbox}[greenbox]{2. Identifier les ramifications}

Tout groupe attaché à la chaîne principale est un substituant.

\begin{itemize}[leftmargin=*]
    \item CH$_3$ : méthyl
    \item CH$_2$CH$_3$ : éthyl
    \item CH(CH$_3$)$_2$ : isopropyl
\end{itemize}

\vspace{0.3cm}

\textbf{Exemple :}

2-méthylhexane

\end{sectionbox}

%\columnbreak

\begin{sectionbox}[purplebox]{3. Groupes caractéristiques courants}

\small

\begin{tabularx}{\linewidth}{|X|X|X|}
\hline
\textbf{Fonction} & \textbf{Préfixe} & \textbf{Suffixe} \\
\hline
Alcool & hydroxy- & -ol \\
\hline
Aldéhyde & formyl- & -al \\
\hline
Cétone & oxo- & -one \\
\hline
Acide carboxylique & carboxy- & -oïque \\
\hline
Amine & amino- & -amine \\
\hline
Alcène & & -ène \\
\hline
Alcyne & & -yne \\
\hline
\end{tabularx}

\end{sectionbox}

%\end{multicols}
%==================== PRIORITE ====================

\begin{sectionbox}[orangebox]{4. Choisir le groupe principal}

Le groupe de plus haute priorité fournit le suffixe du nom.

\vspace{0.3cm}

\begin{center}
\Large
Acide $>$ Ester $>$ Aldéhyde $>$ Cétone $>$ Alcool $>$ Amine $>$ Alcène $>$ Alcyne $>$ Alcane
\end{center}

\end{sectionbox}

\vspace{0.3cm}

%==================== CONSTRUCTION ====================

\begin{sectionbox}[cyanbox]{5. Construire le nom de la molécule}

\begin{center}
\Large
\textbf{Préfixes + chaîne principale + insaturation + suffixe}
\end{center}

\vspace{0.4cm}

\begin{tabularx}{\linewidth}{|X|X|X|X|}
\hline
\textbf{Étape} & \textbf{Exemple} & \textbf{Rôle} & \textbf{Exemple obtenu} \\
\hline
Substituants & 2-méthyl & Préfixe & 2-méthyl \\
\hline
Chaîne & hex- & Nombre de carbones & hex \\
\hline
Insaturation & -ène & Double liaison & hex-2-ène \\
\hline
Fonction & -ol & Groupe principal & hex-2-én-1-ol \\
\hline
\end{tabularx}

\end{sectionbox}

\vspace{0.4cm}

%==================== EXEMPLES ====================

\begin{sectionbox}[lightblue]{Exemples}

\begin{enumerate}
    \item \textbf{2-méthylhexane}
    \item \textbf{3-éthylpent-1-ène}
    \item \textbf{acide 2-méthylbutanoïque}
\end{enumerate}

\end{sectionbox}

\end{document}